\documentclass[10pt,a4paper]{article}
\usepackage{amsmath,amssymb,amsthm,graphicx}
\newcommand{\bits}{\{0,1\}}
\newcommand{\DT}{\operatorname{DT}}
\newcommand{\DTSize}{\operatorname{DTSize}}
\newcommand{\CC}{\operatorname C}
\newcommand{\comp}{\diamond}
\newcommand{\OR}{\operatorname{OR}}
\newcommand{\AND}{\operatorname{AND}}
\newcommand{\XOR}{\operatorname{XOR}}
\newenvironment{originalresult}[1]{\par\medskip\noindent\textbf{#1.}\itshape}{\par\medskip}
\title{Complete gadget trichotomy for parity decision-tree lifting}
\author{Yaroslav Alekseev, Yuval Filmus, Alexander V. Smal}
\date{doi:10.1007/s00037-025-00276-5; CC BY 4.0}
\begin{document}
\maketitle
\noindent\textit{Editorial note.} Actual complete CC BY4.0 journal version. Faithful editable transcription retaining all source-local proofs used by Theorem5.15, including the full OR-certificate lifting argument. Source notation and formulas are preserved, including apparent author typographical choices; no author mathematics is repaired. Source Lemma5.18 and log-rank proofs are not needed for this original certificate-based proof and are omitted. Standard prior affine-closure and certificate relations remain exact cited prerequisites.\par\medskip
\noindent\textit{Allowed standard prerequisites (editorial).}\begin{itemize}
\item Boolean affine-subspace closure characterization. Chen2009 Theorem4.3; exact used source Proposition5.13 retained
\item Polynomial certificate/decision-tree relations, including parity variants. Beals etal.2001 and Zhang/Shi2010; exact used relations retained in final proof
\end{itemize}\medskip
\section*{Block composition (author source p.2)}
For $f:\bits^n\to V$ and $g:\bits^m\to\bits$, a (block-)composition $f\comp g:\bits^{n\times m}\to V$ is defined by
\[(f\comp g)(z_1,z_2,\ldots,z_n):=f(g(z_1),g(z_2),\ldots,g(z_n)),\]
where each $z_i\in\bits^m$.
\section*{Prerequisites (author source pp.10--12)}
All the logarithms are base 2. For $g:\bits^k\to\bits$, we define a function $g^n:\bits^{n\times k}\to\bits^n$ such that
\[g^n(x_1,x_2,\ldots,x_n):=(g(x_1),g(x_2),\ldots,g(x_n)).\]
For a function $f:\bits^n\to\bits$, an input $x\in\bits^n$, and a partial assignment $\sigma\in\{0,1,*\}^n$ that agrees with $x$ on all non-$*$ coordinates, we use $f|_\sigma$ to denote a restriction of $f$ to $\sigma$ and $x|_\sigma$ to denote a projection of $x$ to the $*$-coordinates of $\sigma$.
\paragraph{Decision tree complexity.} A decision tree is a rooted binary tree with internal nodes labeled by input variables (representing queries), edges labeled with 0 and 1, and leaves labeled by some values. Decision tree evaluation is defined in a natural way: given an assignment for the input variables, we traverse the tree starting from the root and then sequentially choose every next edge according to the given assignments until we reach some leaf. The label of this leaf is the result of the evaluation. A decision tree computes some function $f$ if for all possible assignments to the input variables, the decision tree evaluates to the value of the function. For a function $f$, $\DT(f)$ denotes the minimal depth of a decision tree computing $f$, and $\DTSize(f)$ denotes the minimal number of leaves in a decision tree computing $f$.
\paragraph{Certificate complexity.} A certificate complexity is a non-deterministic analogue of the decision tree complexity. A certificate for an input $x\in\bits^n$ to a function $f$ is a set $S\subseteq[n]$ of indices such that $f$ restricted to all inputs that match $x$ on $S$ is constant, i.e., $f(y)=f(x)$ whenever $y|_S=x|_S$. The certificate complexity $\CC(f,x)$ of $f$ at input $x$ is the size of the smallest certificate for $x$. Finally, the certificate complexity of function $f$ is defined as $\CC(f):=\max_{x\in\bits^n}\CC(f,x)$.

\section*{Complexity definitions (author source pp.11--13)}
\paragraph{Parity decision tree complexity.} A parity decision tree extends the concept of a decision tree by allowing queries that are arbitrary linear combinations of the input variables. It can be described as a rooted binary tree with internal nodes labeled by linear combinations of the input variables, edges labeled with 0 and 1, and leaves labeled by some values. The evaluation is defined analogously. For a function $f$, $\DT_\oplus(f)$ denotes the minimal depth of a parity decision tree computing $f$, and $\DTSize_\oplus(f)$ denotes the minimal number of leaves in a parity decision tree computing $f$.

\section*{Parity certificates (author source pp.12,24)}
A parity certificate for an input $x\in\bits^n$ to a function $f$ is a set $A$ of linear forms such that $f$ is constant on the affine subspace defined by $A=A(x)$. A parity certificate complexity $C_\oplus(f,x)$ of a function $f$ at input $x$ is the size (number of linear forms) of the smallest parity certificate for $x$. The parity certificate complexity of $f$ is defined as $C_\oplus(f):=\max_{x\in\bits^n}C_\oplus(f,x)$.
\section*{5.2 Minimum weight lemma}
A parity certificate $C$ is a system of affine equations over input variables. It can be represented by a matrix $A$ and a vector $b$ that define a linear subspace, i.e., $Ax=b$.
\begin{originalresult}{Definition 5.2}
The minimum weight of a parity certificate $C$ is the minimum Hamming weight of nonzero vectors in the row space of $A$.
\end{originalresult}
For example, one can consider a certificate $\{x+y=1,x+y+z=0\}$. The minimum weight of this certificate is equal to 1, and this corresponds to the equation $z=1$. On the other hand, the minimum weight of a certificate $\{x+y=1,x+z=0\}$ is equal to 2.

\begin{originalresult}{Lemma 5.3 (minimal weight lemma; author source p.25)}
For $f:\bits^n\to\bits$ and $g:\bits^m\to\bits$, suppose that $f\comp g$ has parity certificate complexity at most $k$ at some point $x$. Let $K$ be a parameter. If the certificate complexity of $f$ at $g^n(x)$ is greater than $k\cdot K$ then we can find a blockwise partial assignment $\sigma$ to the inputs of $f\comp g$ consistent with $x$, such that for some $k'\le k$:
\begin{itemize}
\item $(f\comp g)|_\sigma$ has a parity certificate of size $k'$ at $x|_\sigma$ whose minimum weight is at least $K$,
\item $C(f|_{g^n(\sigma)},g^n(x)|_{g^n(\sigma)})>k'K$. (Note that $\sigma$ is blockwise and agrees with $x$, so here we consider the certificate complexity of $f$ at $g^n(x)$ on all the blocks that are not set by $\sigma$).
\end{itemize}
\end{originalresult}
\begin{proof}
The proof is by induction on $k$. If $k=0$, then the premises contradict each other, so we can assume that $k>0$.
Let $C$ be the parity certificate for $f\comp g$ at input $x$. By the assumption, the size of $C$ is at most $k$. If its minimum weight is at least $K$, then we are done. Otherwise, without loss of generality, $C$ contains a nontrivial constraint $\ell=\alpha$ whose support is less than $K$. Choose a partial assignment that agrees with $x$ on the support of this constraint, and complete it to a blockwise partial assignment $\sigma$ consistent with $x$ (complete only the blocks with at least one variable assigned). Also, let $\tau:=g^n(\sigma)$.
By construction, $(f\comp g)|_\sigma$ has a parity certificate of size at most $k-1$ at $x|_\sigma$, namely $C\setminus\{\ell=\alpha\}$. On the other hand, since we only set the values of at most $K$ blocks, the certificate complexity of $f$ reduces by at most $K$ after being restricted to $\tau$, and hence $C(f|_\tau,g^n(x)|_\tau)>(k-1)K$. Now we can apply the induction hypothesis for $k-1$.
\end{proof}
The following lemma demonstrates an essential property of the parity certificates with huge minimal weight:
\begin{originalresult}{Lemma 5.4 (author source pp.25--26)}
Suppose that $C_\oplus$ is a parity certificate with minimum weight at least $K$ and $C$ is a (regular) certificate of size less than $K$. Then there is a point that satisfies both $C$ and $C_\oplus$.
\end{originalresult}
In the proof, the parity certificate is processed with a careful Gaussian elimination, and then the variables fixed by the certificate $C$ are free variables of the parity certificate.
\begin{proof}
Let $V+b$ be the affine subspace corresponding to $C_\oplus$, where $V$ is a linear subspace. We pick a basis $\ell_1,\ldots,\ell_k$ for $V$ and indices $i_1,\ldots,i_k\in[n]$ such that:
\begin{itemize}
\item $z_{i_t}$ is in the support of $\ell_t$,
\item $z_{i_s}$ is not in the support of $\ell_t$ for $s<t$,
\item $z_{i_t}$ does not belong to the support of $C$.
\end{itemize}
We do it inductively. At step $t$, we have already picked $\ell_1,\ldots,\ell_{t-1}$. Consider vectors in $V$ whose support does not contain any of the variables $z_{i_1},\ldots,z_{i_{t-1}}$. Since $V$ has minimum weight $K$, we can find such a vector $\ell_t$ whose support contains a variable $z_{i_t}$ with $z_{i_t}\notin C$. By the assumption, this vector does not contain any of the variables $z_{i_1},\ldots,z_{i_{t-1}}$.
Now we can construct an input $\alpha$ which satisfies both $C$ and $C_\oplus$. First, we set the variables from $C$ such that the resulting input satisfies $C$. Then we set arbitrarily all the variables not occurring in $C$ or $\{z_{i_1},\ldots,z_{i_k}\}$. Finally, we set variables $z_{i_k},\ldots,z_{i_1}$ (in this order) in such a way that the resulting input satisfies $C_\oplus$.
\end{proof}

\section*{5.3.2 Lifting (author source pp.27--30)}
Given Lemma 5.3, we are going to prove the following lifting for the certificate complexity.
\begin{originalresult}{Theorem 5.7}
For every function $f:\bits^n\to\bits$, we have $C(f)=O(C_\oplus(f\comp\OR)^3)$.
\end{originalresult}
Throughout the proof of the theorem, we will think of $f$ as having inputs $x_1,x_2,\ldots,x_n$, and of block composition $f\comp\OR$ as having inputs $z_{1,1},z_{1,2},z_{2,1},z_{2,2},\ldots,z_{n,1},z_{n,2}$, such that $x_i$ corresponds to $z_{i,1}\vee z_{i,2}$. The twin of $z_{i,1}$ is $z_{i,2}$, and vice versa. We start the proof with the following lemma, another technical tool in this section. Together with Lemma 5.3, the proof of the lifting theorem is essentially a combination of these two tricks.
\begin{originalresult}{Lemma 5.8}
Suppose that $C$ is a parity certificate for $f\comp\OR$ of size $k$ whose minimum weight is at least $k$. Then there are $k$ indices $i_1,\ldots,i_k$ such that for every input $x$ with $x_{i_1}=\cdots=x_{i_k}=1$, we can find an input $z$ satisfying $C$ such that $\OR^n(z)=x$. In particular, there is an input $z$ satisfying $C$ such that $\OR^n(z)=(1,\ldots,1)$.
\end{originalresult}
\begin{proof}
Let $V+b$ be the affine subspace corresponding to $C$, where $V$ is a linear subspace (the row space of the linear system $Ax=b$ in Definition 5.2). We pick a basis $\ell_1,\ldots,\ell_k$ for $V$ and indices $(i_1,j_1),\ldots,(i_k,j_k)\in[n]\times[2]$ such that:
\begin{itemize}
\item $z_{i_t,j_t}$ is in the support of $\ell_t$.
\item $z_{i_s,1},z_{i_s,2}$ are not in the support of $\ell_t$ for $s<t$.
\end{itemize}
We build the basis by successively adding vectors to it, similar to the Gaussian elimination. At step $t$, we have already chosen basis vectors $\ell_1,\ldots,\ell_{t-1}$ and the corresponding indices. Consider all the vectors in $V$ whose support does not contain any of the variables $z_{i_1,j_1},\ldots,z_{i_{t-1},j_{t-1}}$. Since $V$ has minimum weight $k>t-1$, we can find a vector $\ell_t$ whose support contains a variable $z_{i_t,j_t}$ with $i_t\notin\{i_1,\ldots,i_{t-1}\}$ and add it to the basis.
Now let $x$ be an input satisfying $x_{i_1}=\cdots=x_{i_k}=1$. We set the twins of $z_{i_1,j_1},\ldots,z_{i_t,j_t}$ to 1, and set $z_{i,1},z_{i,2}$ for $i\ne i_1,\ldots,i_t$ in an arbitrary way subject to $z_{i,1}\vee z_{i,2}=x_i$. We then set $z_{i_t,j_t},\ldots,z_{i_1,j_1}$ (in this order) in such a way that the resulting input satisfies $C$.
\end{proof}
The next lemma that we need for the proof of the lifting result gives an upper bound on the minimum weight of a parity certificate in terms of the sizes of certificates for 0 and 1.
\begin{originalresult}{Lemma 5.9}
Suppose that $f'\comp\OR$ has a parity certificate $C'$ of size $k'$ which certifies $b$, and a parity certificate $C''$ of size $k''$ which certifies $\bar b$. Then the minimum weight of $C'$ is at most $k'+2(k'')^2$.
\end{originalresult}
\begin{proof}
First, observe that both $k'>0$ and $k''>0$ since otherwise $f'\comp\OR$ is constant and cannot have certificates for both 0 and 1.
Consider the case $k''=1$. The certificate $C''$ consists of one constraint $\ell=\alpha$. By Lemma 5.8, since $C''$ has size 1, some input $z$ satisfying $C''$ satisfies $\OR^n(z)=(1,\ldots,1)$, and so $f'(1,\ldots,1)=\bar b$. Suppose for the sake of contradiction that the minimum weight of $C'$ is at least $k'+2(k'')^2+1=k'+3>k'$. According to Lemma 5.8, some input $z$ satisfying $C'$ satisfies also $\OR^n(z)=(1,\ldots,1)$, and so $f'(1,\ldots,1)=b$, which is impossible. So the minimum weight of $C'$ is at most $k'+2$.
We proceed by induction on $k''$. Suppose for the sake of contradiction that the minimum weight of $C'$ is at least $k'+2(k'')^2+1>k'$. According to Lemma 5.8, some input $z$ satisfying $C'$ satisfies $\OR^n(z)=(1,\ldots,1)$, and so $f'(1,\ldots,1)=b$. If the minimum weight of $C''$ was also at least $k''$, then the same argument would imply that $f'(1,\ldots,1)=\bar b$. Therefore, the minimum weight of $C'$ is less than $k''$. Consider some constraint $\ell=\alpha$ of $C''$ of weight at most $k''$. Lets choose some partial assignment satisfying $\ell=\alpha$, and complete it to a blockwise partial assignment $\tau$ in an arbitrary way. Thus $\tau$ assigns at most $2k''$ variables. By Lemma 5.4, the assignment $\tau$ cannot contradict $C'$ since the minimum weight of $C'$ is greater than $k'+2(k'')^2>2k''$. Therefore $(f'\comp\OR)|_\tau$ still has a $b$-certificate $C'|_\tau$ of size $k'$ and a $\bar b$-certificate $C''|_\tau$ of size $k''-1$. By induction, we know that the minimum weight of $C'|_\tau$ is at most $k'+2(k''-1)^2$, and so the minimum weight of $C'$ is at most $k'+2(k''-1)^2+2k''<k'+2(k'')^2$.
\end{proof}
The following lemma applies the `minimal weight lemma' to the current setting.
\begin{originalresult}{Lemma 5.10}
Suppose that $f\comp\OR$ has a parity certificate at some input $z$ of size at most $k$. Let $K$ be a parameter. Suppose that the certificate complexity of $f$ at $\OR^n(z)$ is greater than $k\cdot K$. Then we can find a blockwise partial assignment $\sigma$ to the inputs of $f\comp\OR$, consistent with $z$, such that for some $k'\le k$:
\begin{itemize}
\item $(f\comp\OR)|_\sigma$ has a parity certificate of size $k'$ at $z|_\sigma$ whose minimum weight is at least $K$.
\item $C(f|_{\OR^n(\sigma)},\OR^n(z)|_{\OR^n(\sigma)})>k'K$.
\end{itemize}
\end{originalresult}
\begin{proof}
We apply Lemma 5.3 for a function $f$ and a gadget $g=\OR$.
\end{proof}
Now we have all the ingredients for proving Theorem 5.7.
\begin{proof}[Proof of Theorem 5.7]
Let $C_\oplus(f\comp\OR)=k$, and let $K$ be a parameter to be determined later. Suppose that $C(f)>kK$. Apply Lemma 5.10 to get a non-constant function $f'$, obtained by restricting $f$, such that $f'\comp\OR$ has a parity certificate of size $k'$ whose minimum weight is at least $K$ (the function $f'$ is non-constant since its certificate complexity is positive). Since $C_\oplus(f'\comp\OR)\le C_\oplus(f\comp\OR)$ and $f'$ is non-constant, Lemma 5.9 (with $k''\le k$) shows that $K\le k'+2(k'')^2\le k+2k^2$. By choosing $K=k+2k^2+1$ we reach a contradiction. Therefore $C(f)\le kK=O(k^3)$.
\end{proof}

\section*{5.4 AND/OR gadgets (author source pp.30--31)}
A function $g:\bits^m\to\bits$ affine projects to a function $h:\bits^p\to\bits$ if there are affine functions $\ell_1,\ldots,\ell_m:\mathbb Z_2^p\to\mathbb Z_2$ such that
\[g(\ell_1(z),\ldots,\ell_m(z))=h(z).\]
A gadget $g:\bits^m\to\bits$ is AND/OR if it affine projects to both binary AND and binary OR. Here are some examples of AND/OR gadgets:
\begin{itemize}
\item $g(x,y,z):=x\vee(y\wedge z)$. The projections: $g(0,a,b)=a\wedge b$ and $g(a,b,1)=a\vee b$.
\item $g(x,y,z):=[x+y+z=1]$. The projections: $g(1,\bar a,\bar b)=a\wedge b$ and $g(\bar a,\bar b,a\oplus b)=a\vee b$.
\end{itemize}
\begin{originalresult}{Theorem 5.12}
If $g$ is an AND/OR gadget then for all functions $f:\bits^n\to\bits$
\[\log\DTSize_\oplus(f\comp g)\ge C(f).\]
\end{originalresult}
\begin{proof}
Let $T_0$ be a parity decision tree for $f\comp g$. Let $z$ be a point such that $C(f,z)=C(f)$. Since $g$ is AND/OR, we can apply affine substitutions to $T_0$ and get a parity decision tree $T$ computing function $f\comp(h_{z_1},\ldots,h_{z_n})$, where $h_0(a,b):=a\wedge b$ and $h_1(a,b):=a\vee b$.
Now we follow the simulation proof of Chattopadhyay et al. (2023). We start at the root. When at a node $v$, if the query is constant then we follow the corresponding edge, otherwise we pick some variable, without loss of generality $a_i$, which appears in the query. So the query is of the form $a_i\oplus\ell$. We set $b_i:=0$ if $z_i=0$, and we set $b_i:=1$ if $z_i=1$. Then we choose $c$, the answer to the query, such that it leads to the smaller subtree, and set $a_i:=\ell\oplus c$. From now on, we see neither $a_i$ nor $b_i$.
Eventually, we reach a leaf. Suppose that this happens after eliminating coordinates $i_1,\ldots,i_K$, which means that $T$ contains at least $2^K$ leaves. For every input $z'$ to $f$ which agrees with $z$ on $i_1,\ldots,i_K$ we can find an input $w'$ to $T$ which satisfies $g(w')=z'$ and leads to the same leaf, showing that $z'_{i_1}=z_{i_1},\ldots,z'_{i_K}=z_{i_K}$ is a certificate for $f$ at $z$. We conclude that $K\ge C(f,z)=C(f)$, and so $|T_0|\ge|T|\ge2^K$.
\end{proof}
\section*{5.5 Gadget classification (author source pp.31--33)}
For the proof, we will need the following classical fact.
\begin{originalresult}{Proposition 5.13 (Chen2009, Theorem4.3)}
Suppose that $f:\bits^n\to\bits$ satisfies the following condition: whenever $f(x)=f(y)=f(z)=1$ then $f(x\oplus y\oplus z)=1$. Then either $f\equiv0$, or $f$ is an indicator of an affine subspace.
\end{originalresult}
We need this lemma to classify gadgets that are not AND/OR.
\begin{originalresult}{Lemma 5.14}
If $g$ is not an AND/OR then $g$ is a disjunction or a conjunction of affine forms in the inputs.
\end{originalresult}
\begin{proof}
Suppose that $g$ does not affine project to OR. We claim that for any inputs $x,y,z$, if $g(x)=g(y)=g(z)=1$ then $g(x\oplus y\oplus z)=1$. Indeed, suppose that $g(x)=g(y)=g(z)=1$ but $g(x\oplus y\oplus z)=0$, and consider the projection
\[h(a,b):=g(x\oplus y\oplus z\oplus a(x\oplus y)\oplus b(x\oplus z)).\]
Then $h(0,0)=g(x\oplus y\oplus z)=0$, $h(1,0)=g(z)=1$, $h(0,1)=g(y)=1$, and $h(1,1)=g(x)=1$. Thus, we can conclude that $g$ affine projects to OR.
Applying Proposition 5.13, we see that if $g$ does not affine project to OR then it is either the empty disjunction or a conjunction of affine forms. Similarly, if $g$ does not affine project to AND then $\bar g$ does not affine project to OR, and so $g$ is either the empty conjunction or a disjunction of affine forms.
\end{proof}
Now we can prove the classification of gadgets.

\begin{originalresult}{Theorem 5.15}
For every gadget $g$, one of the following cases holds:
\begin{enumerate}
\item[(i)] There is an infinite family of functions $f_n$ with $\DT(f_n)\to\infty$ and $\DTSize_\oplus(f_n\comp g)=O(1)$.
\item[(ii)] For every function $f$, $\DT_\oplus(f\comp g)=\Omega(\DT(f)^{\Omega(1)})$. There is an infinite family of functions $f_n$ with $\DT(f_n)\to\infty$ and $\DTSize(f_n\comp g)=O(\DT(f))$.
\item[(iii)] For every function $f$, $\log\DTSize_\oplus(f\comp g)=\Omega(\DT(f)^{\Omega(1)})$ and also $C_\oplus(f\comp g)=\Omega(C(f)^{\Omega(1)})$.
\end{enumerate}
\end{originalresult}

\begin{proof}
If gadget $g$ is AND/OR then by Theorem 5.12 we have
\[\log\DTSize_\oplus(f\comp g)\ge C(f)=\Omega(\sqrt{\DT(f)}).\]
Since $C_\oplus(f\comp g)=\Omega(\sqrt{\DT_\oplus(f\comp g)})$ (Zhang \& Shi2010) and also $\DT_\oplus(f\comp g)\ge\log\DTSize_\oplus(f\comp g)$, we get the third case of the theorem.
Otherwise, by Lemma 5.14, gadget $g$ is either a conjunction of affine forms or a disjunction of affine forms. Suppose, without loss of generality, that $g$ is a disjunction of affine forms. If $g\equiv0$ or $g$ is an affine form, then we can take $f_n=\oplus_n$ to get the first case of the theorem. If $g$ is not affine then it affine projects to OR, and thus we get the second case: we need to apply Theorem 5.7 together with the polynomial relations between (parity) certificate complexity and (parity) decision tree complexity; the infinite family of functions is $f_n=\OR_n$.
\end{proof}

\begin{thebibliography}{9}
\bibitem{beals} Robert Beals, Harry Buhrman, Richard Cleve, Michele Mosca and Ronald de Wolf (2001). Quantum Lower Bounds by Polynomials. J.ACM48(4),778--797. doi:10.1145/502090.502097.
\bibitem{chen} Hubie Chen (2009). A Rendezvous of Logic, Complexity, and Algebra. ACM Comput.Surv.42(1). doi:10.1145/1592451.1592453.
\bibitem{stifling} Arkadev Chattopadhyay, Nikhil S.Mande, Swagato Sanyal and Suhail Sherif (2023). Lifting to parity decision trees via stifling. arXiv:2211.17214.
\bibitem{zhangshi} Zhiqiang Zhang and Yaoyun Shi (2010). On the parity complexity measures of Boolean functions. Theoretical Computer Science411(26),2612--2618.
\end{thebibliography}
\end{document}
